%#!latexmk -pdfdvi -f main.tex
%#LPR acroread %s

\RequirePackage{plautopatch}
\documentclass[paper=a4,
fontsize=12pt,
reqno,
uplatex,
dvipdfmx]{jlreq}
%
\jlreqsetup{
thebibliography_heading={\section*{\refname}},
}

%
%%%%%%%%%%%%%%% Include Package file %%%%%%%%%%%%%%%
\usepackage{amsmath,amssymb,amsthm}
%\usepackage{mathrsfs}
\usepackage{graphicx,color}
\usepackage{newtxtext}
\usepackage[varg]{newtxmath}
\usepackage[shortlabels]{enumitem}
\usepackage[haranoaji]{pxchfon}

\usepackage[alphabetic, y2k]{amsrefs}
%
%%%%%%%%%%%%%%% Paper size settings %%%%%%%%%%%%%%%
%
%
\usepackage{geometry}%
\geometry{left=20mm,right=20mm,top=15mm,bottom=25mm}
\geometry{a4paper}
%

%%%%%%%%%%%%%%% mathtools %%%%%%%%%%%%%%%
%
\usepackage{mathtools}
\mathtoolsset{showonlyrefs=true}
%
%
%
%%%%%%%%%%%%%%% hyperref %%%%%%%%%%%%%%%
%
\usepackage{hyperref}
\usepackage{pxjahyper}
%
%
%
%%%%%%%%%%%%%%% Setting of Theorem-like environment %%%%%%%%%%%%%%%
% [Italic font in Theorem-like environment]
%\theoremstyle{definition} % if change the roman font
\theoremstyle{definition} % Change the Roman font
%\newtheorem{definition}{定義}[chapter]
\newtheorem{theorem}{定理}
\newtheorem{proposition}[theorem]{命題}
\newtheorem{lemma}[theorem]{補題}
\newtheorem{corollary}[theorem]{系}
\newtheorem{problem}{問題}
\newtheorem{example}{例}
\newtheorem*{definition}{定義}
\newtheorem*{remark}{注意}
%
 
%
%%%%%%%%%% Proof %%%%%
\renewcommand{\proofname}{証明}  
% ((Use the environment in proof))
% \begin{proof}
% \end{proof}
% white box in the end of proof
%
%
% make new QED
% \QED{foo}
% we write a black box in the right and output foo after the black box
\newcommand{\QED}[1]{\hfill$\blacksquare$ #1\par}
% 
%
%
%%%%%%%%% 1.1 times Line stretch %%%%%%%%%
% If you do not use Japanese, that command should be commented out.
%
%
%%%%%%% Mathematical symbols %%%%%%%%%%%%%
%
\input{mymacro.tex}
%
%%%%%%%%%%%%%%% Counter setting in the equation's reference %%%%%%%%%%%%%%%
% For amsart, use this counter setting
%
%
%%% Definition of the \section (not use the \scshape)
%
%
%\usepackage{input/myamssec}
%
%%% \par after proofname
%
%\usepackage{input/myamspf}
%
%%% \par after theorem environment
%
%\usepackage{input/myamsthm}
%

%%%%%%%%%% Framed settings %%%%%
%
%\definecolor{shadecolor}{gray}{0.9}
%\newenvironment{shadedf}{%
%  \def\FrameCommand{\fcolorbox{white}{shadecolor}}%
%  \MakeFramed {\advance\hsize-\width \FrameRestore}%
%  }%
%{\endMakeFramed}
%
%
%


%%%%%%%%%%%%%%% temporary Package %%%%%%%%%%%%%%%

\usepackage[inline,right]{showlabels}   % Output of cite and reference

%\usepackage{showkeys}   % Output of cite and reference
%\usepackage{refcheck}   % check of no uselabel 

%%%%%%%%%%%%%%% title Setting %%%%%%%%%%%%%%%
%
\title{Hilbert積分の拡張と応用}
\author{}
%
%
%%%%%%%%%%%%%%% page style setting %%%%%%%%%%%%%%%
%
\allowdisplaybreaks[1]
\DeclareMathOperator{\rot}{rot}
\let\Re\relax
\let\Im\relax
\DeclareMathOperator{\Re}{Re}
\DeclareMathOperator{\Im}{Im}

\usepackage{bm}
\renewcommand{\vec}[1]{\bm{#1}}


%
%%%%%%%%%%%%%%% end preamble %%%%%%%%%%%%%%%
%
\begin{document}
\maketitle

\section{Hilbert積分}
\cite{MR944909}*{Chapter 9}によると，Hilbertは2重和
\begin{equation}
 \sum_{m=1}^\infty
  \sum_{n=1}^\infty
  \frac{1}{m+n}a_mb_n
\end{equation}
の研究を始めて行ったとされる．この積分型
\begin{equation}
 \int_0^\infty
  \left(
  \int_0^\infty
  \frac{1}{x+y}
  f(x)
  g(y)
  \,dx
  \right)
  \,dy
\end{equation}
の一般化である
\begin{equation}
 \int_0^\infty
  \left(
  \int_0^\infty
  K(x,y)
  f(x)
  g(y)
  \,dx
  \right)
  \,dy
\end{equation}
を考えよう．

$K:(0,\infty)\times(0,\infty)\rightarrow\R$が$-1$次斉次であるとはすべての$x,y>0$と$\lambda>0$に対して
\begin{equation}
 K(\lambda x,\lambda y)
  =
  \frac{1}{\lambda}K(x,y)
  =
  \lambda^{-1}K(x,y)
\end{equation}
をみたすことである．

\begin{proposition}
 $K:(0,\infty)\times(0,\infty)\rightarrow\R$を
 $(0,\infty)\times(0,\infty)$上の非負値，$-1$次斉次なLebesgue可測関数
 とする．このとき，任意の$f,g\in C_0^\infty(0,\infty)$と$0<a\leq\infty$，
 $1<p<\infty$に対して
 \begin{equation}
  \label{eq:HilbertIntegral1}
    \int_0^a
    \left(
     \int_0^\infty
     K(x,y)
     |g(y)|
     dy
    \right)
    |f(x)|
    \,dx
   \leq
   \int_0^\infty
   K(1,w)
   w^{-\frac{1}{q}}
   \|g\|_{L^{q}(0,wa)}
   \,dw
   \|f\|_{L^p(0,a)}
 \end{equation}
 が成り立つ．ただし，$q$は$p$のH\"older共役，すなわち
 $\frac{1}{p}+\frac{1}{q}=1$である．
\end{proposition}

\begin{proof}
 $0<x<a$に対して，$y=wx$と変数変換すると，$K$が$-1$次斉次より
 \begin{equation}
  \int_0^\infty
   K(x,y)|g(y)|
   \,dy
   =
   \int_0^\infty
   K(x,wx)|g(wx)|
   x
   \,dw
   =
   \int_0^\infty
   K(1,w)|g(wx)|
   \,dw
 \end{equation}
 が成り立つ．Fubiniの定理より
 \begin{equation}
  \label{eq:HilbertIntegral_proof}
  \int_0^a
   \left(
     \int_0^\infty
   K(x,y)|g(y)|
   \,dy
   \right)
   |f(x)|
   \,dx
   =
   \int_0^\infty
   K(1,w)
   \left(
    \int_0^a
    |f(x)|
    |g(wx)|
    \,dx
   \right)
   \,dw
 \end{equation}
 となる．H\"olderの不等式と変数変換$y=wx$より
 \begin{equation}
  \begin{split}
   \int_0^a
   |f(x)|
   |g(wx)|
   \,dx
   &\leq
   \|f\|_{L^p(0,a)}
   \left(
   \int_0^a
   |g(wx)|^q
   \,dx
   \right)^{\frac{1}{q}}
   \\
   &=
   \|f\|_{L^p(0,a)}
   \left(
   \int_0^{wa}
   |g(y)|^q
   w^{-1}\,dy
   \right)^{\frac{1}{q}}
   =
   \|f\|_{L^p(0,a)}
   \|g\|_{L^p(0,wa)}
   w^{-\frac{1}{q}}
  \end{split} 
 \end{equation}
 となるから
 \begin{equation}
  \int_0^\infty
   K(1,w)
   \left(
    \int_0^a
    |f(x)|
    |g(wx)|
    \,dx
   \right)
   \,dw
   \leq
   \int_0^\infty
   K(1,w)
   \|g\|_{L^p(0,wa)}
   w^{-\frac{1}{q}}
   \,dw
   \|f\|_{L^p(0,a)}
 \end{equation}
 が得られる．~\eqref{eq:HilbertIntegral_proof}に代入すること
 で~\eqref{eq:HilbertIntegral1}が得られる．
\end{proof}

双対性を用いることで次が得られる．

\begin{corollary}
 $K:(0,\infty)\times(0,\infty)\rightarrow\R$を
 $(0,\infty)\times(0,\infty)$上の非負値，$-1$次斉次なLebesgue可測関数
 とする．このとき，任意の$g\in C_0^\infty(0,\infty)$と$0<a\leq\infty$，
 $1<p<\infty$に対して
 \begin{equation}
  \label{eq:HilbertIntegral2}
  \int_0^a
   \left(
   \int_0^\infty
   K(x,y)g(y)
   \,dy
   \right)^p
   \,dx
   \leq
   \left(
   \int_0^\infty
   K(1,w)
   w^{-\frac{1}{p}}
   \|g\|_{L^p(0,wa)}
   \,dw
   \right)^p
 \end{equation}
 が成り立つ．
\end{corollary}

\begin{proof}
 任意の$\phi\in C_0^\infty(0,\infty)$に対して~\eqref{eq:HilbertIntegral1}
 より
 \begin{equation}
  \begin{split}
   \left|
   \int_0^a
   \left(
   \int_0^\infty
   K(x,y)
   |g(y)|
   dy
   \right)
   \phi(x)
   \,dx
   \right|
   &\leq
   \int_0^a
    \left(
     \int_0^\infty
     K(x,y)
     |g(y)|
     dy
    \right)
    |\phi(x)|
   \,dx
   \\
   &\leq
   \int_0^\infty
   K(1,w)
   w^{-\frac{1}{q}}
   \|g\|_{L^{q}(0,wa)}
   \,dw
   \|\phi\|_{L^p(0,a)}
  \end{split} 
 \end{equation}
 が成り立つ．双対性の議論により~\eqref{eq:HilbertIntegral2}が得られる．
\end{proof}

\section{Hilbert積分の応用}

$\beta>0$と$x,y>0$に対して
\begin{equation}
 K(x,y)
  :=
  \left(
  \frac{y}{x}
  \right)^\beta
  \frac{1}{y}
  \chi_{(0,x)}(y)
\end{equation}
とおく．任意の$\lambda>0$に対して
\begin{equation}
 K(\lambda x,\lambda y)
  =
  \left(
  \frac{y}{x}
  \right)^\beta
  \frac{1}{\lambda y}
  \chi_{(0,\lambda x)}(\lambda y)
  =
  \frac{1}{\lambda}
  K(x, y)
\end{equation}
となるから，$K$は$-1$次斉次な非負値Lebesgue可測関数である．従っ
て~\eqref{eq:HilbertIntegral2}より
\begin{equation}
 \label{eq:Hardy01}
\begin{split}
  \int_0^a
  \left(
  \int_0^x
  \left(
  \frac{y}{x}
  \right)^\beta
  |g(y)|
  \frac{dy}{y}
  \right)^p
  \,dx
  &\leq
  \left(
  \int_0^1
  w^{\beta-1-\frac{1}{p}}
  \left(
  \int_0^{wa}
  |g(y)|^p
  \,dy
  \right)^\frac{1}{p}
  \,dw
  \right)^p
  \\
  &\leq
  \left(
  \int_0^1
  w^{\beta-1-\frac{1}{p}}
  \,dw
  \right)^p
  \int_0^{a}
  |g(y)|^p
  \,dy
 \end{split}
\end{equation}
が得られる．$\phi(y):=y^\beta|g(y)|$とおくと
\begin{equation}
 \text{~\eqref{eq:Hardy01}の左辺}
  =
  \int_0^a
  x^{-p\beta+1}
  \left(
    \int_0^x
  \phi(y)
  \,\frac{dy}{y}
  \right)^p
  \,\frac{dx}{x}
\end{equation}
と
\begin{equation}
 \text{~\eqref{eq:Hardy01}の右辺}
  =
  \left(
  \int_0^1
  w^{\beta-1-\frac{1}{p}}
  \,dw
  \right)^p
  \int_0^{a}
  y^{-p\beta+1}\phi(y)^p
  \,\frac{dy}{y}
\end{equation}
より$\alpha>0$に対して
\begin{equation}
 -p\beta+1+-p\alpha
  \Longleftrightarrow
  \beta
  =
  \frac{1}{p}+\alpha
\end{equation}
とおくと
\begin{equation}
 \int_0^a
  x^{-p\alpha}
  \left(
   \int_0^x
  \phi(y)
  \,\frac{dy}{y}
  \right)^p
  \,\frac{dx}{x}
  \leq
  \left(
  \int_0^1
  w^{\alpha-1}
  \,dw
  \right)^p
  \int_0^{a}
  y^{-p\alpha}\phi(y)^p
  \,\frac{dy}{y}
  =
  \frac{1}{\alpha^p}
  \int_0^{a}
  y^{-p\alpha}\phi(y)^p
  \,\frac{dy}{y}
\end{equation}
となる．

\begin{proposition}[\cite{MR2523200}*{(A.10)}]
 $0<a\leq\infty$に対し$\phi:(0,a)\rightarrow\R$を非負値Lebesgue可測関数
 とする．このとき，任意の$\alpha>0$と$p>1$に対して
 \begin{equation}
  \label{eq:Hardy_type_Inequality}
   \int_0^a
   t^{-p\alpha}
   \left(
    \int_0^t
    \phi(s)
    \,\frac{ds}{s}
   \right)^p
   \,\frac{dt}{t}
  \leq
  \frac{1}{\alpha^p}
  \int_0^{a}
  s^{-p\alpha}\phi(s)^p
  \,\frac{ds}{s}   
 \end{equation}
 が成り立つ．
\end{proposition}

\begin{remark}
 Fubiniの定理を用いれば\eqref{eq:Hardy_type_Inequality}は$p=1$でも成立す
 る．
\end{remark}

$X$をBanach空間とする．$1\geq p>\infty$と区間$I\subset\R$に対して，$f\in
L^p_*(I:X)$であるとは，$f:I\rightarrow X$がBochner可積分であって
\begin{equation}
 \|f\|_{L^p_*(I:X)}
  :=
  \left(
  \int_I
  \|f(t)\|_X^p
  \,\frac{dt}{t}
  \right)^{\frac{1}{p}}
  <\infty
\end{equation}
をみたすことである．~\eqref{eq:Hardy_type_Inequality}を用いると，次が成
り立つ．

\begin{corollary}[\cite{MR2523200}*{Corollary A.13}]
 $X$をBanach空間とする．$0<a\leq \infty$，$0<\theta<1$，$1\leq p<\infty$
 に対し，$u:(0,a)\rightarrow X$はBochoner可積分関数であって
\begin{equation}
 t\mapsto u_\theta(t)=t^\theta u(t)
\end{equation}
 が$L^p(0,a:X)$に属するとする．このとき，$t>0$に対して積分平均
 \begin{equation}
  v(t)
   :=
   \frac{1}{t}
   \int_0^tu(t)\,dt
 \end{equation}
 も同様の性質を持ち
 \begin{equation}
  \label{eq:Bochoner-Mean-Inequality}
   \|v_\theta\|_{L^p_*(0,a:X)}
   \leq
   \frac{1}{1-\theta}
  \|u_\theta\|_{L^p_*(0,a:X)}
 \end{equation}
 が成り立つ．
\end{corollary}

\begin{proof}
 $v_\theta$の定義より
 \begin{equation}
  \|v_\theta\|_{L^p_*(0,a:X)}^p
   =
   \int_0^a
   \|v_\theta(t)\|_X^p
   \,\frac{dt}{t}
   =
   \int_0^a
   \|t^\theta v(t)\|_X^p
   \,\frac{dt}{t}
   =
   \int_0^a
   t^{p\theta}
   \|v(t)\|_X^p
   \,\frac{dt}{t}
 \end{equation}
 である．$v$の定義と三角不等式より
 \begin{equation}
  \|v(t)\|_X
   \leq
   t^{-1}
   \int_0^t
   \|u(s)\|_X
   \,ds
   =
  t^{-1}
  \int_0^t
  \|su(s)\|_X
  \,\frac{ds}{s}
 \end{equation}
 だから~\eqref{eq:Hardy_type_Inequality}より
 \begin{equation}
  \begin{split}
   \int_0^a
   t^{p\theta}
   \|v(t)\|_X^p
   \,\frac{dt}{t}
   &\leq
   \int_0^a
   t^{-p(1-\theta)}
   \left(
   \int_0^t
   \|su(s)\|_X
   \,\frac{ds}{s}
   \right)^p
   \,\frac{dt}{t}
   \\
   &\leq
   \frac{1}{(1-\theta)^p}
   \int_0^a
   s^{-p(1-\theta)}
   \|su(s)\|^p_X
   \,\frac{ds}{s}
   \\
   &=
   \frac{1}{(1-\theta)^p}
   \int_0^a
   \|s^\theta u(s)\|^p_X
   \,\frac{ds}{s}
   \\
   &=
   \frac{1}{(1-\theta)^p}
   \int_0^a
   \|u_\theta (s)\|^p_X
   \,\frac{ds}{s}
  \end{split} 
 \end{equation}
 となる．従って，~\eqref{eq:Bochoner-Mean-Inequality}が成り立つ．
\end{proof}



% bibtex Settings
\bibliography{references}
%\bibliographystyle{myjalpha} % in Japanese

%\bibliographystyle{amsplain} % not in Japanese


\end{document}

